\section{Conclusion}\label{sec:conclusion}
%==============================================================================

Deterministic partial views induce a confusability graph on latent tuples. In the full coordinate-view model, the realizable confusability relations are exactly the upward-closed agreement-set families; exact recovery with a $T$-ary auxiliary tag is equivalent to $T$-colorability; block composition lifts the one-shot graph to strong powers; and the normalized block rates converge to the Shannon capacity of the induced graph. The same graph carries the standard Lov\'asz-$\vartheta$ upper theory, and transitivity of confusability gives a model-side route to the cluster-graph equality case.\leanmeta{\LH{MFT3}, \LH{MFT20}, \LH{MFT35}, \LH{MFT48}, \LH{MFT69}, \LH{MFT96}, \LH{MFT102}, \LHrng{MFT}{118}{131}}

Under an affine restriction on the realized state family, the coordinate side carries a representable matroid whose rank gives tractable upper bounds on one-shot and asymptotic recoverability. The same framework also yields a boundary criterion for verifiable zero-error realizations: unit rate is the unique zero-incoherence regime, and achieving it for structural facts requires both zero-delay synchronization and structural side-information. Appendix~\ref{sec:appendix-classification} records representative model instantiations of that criterion.\leanmeta{\LHrng{AFM}{1}{14}, \LHrng{COH}{1}{2}, \LH{RAT1}, \LH{SOT1}}

Several questions remain open. The paper does not classify which upward-closed agreement-set families yield transitive confusability, nor the model-generated graph class up to unlabeled isomorphism. The restricted-state and support-overlap extensions are formalized but not developed textually, and the sharpness of the Lovász-$\vartheta$ upper theory on the non-clique subclass is unresolved. The Lean~4 formalization and submission source bundle are archived at Zenodo~\cite{paper2_partial_views_artifact}.

\subsection*{Declaration of Generative AI and AI-Assisted Technologies in the Manuscript Preparation Process}

Generative AI tools (including Codex, Claude Code, Augment, Kilo, and OpenCode) were used throughout this manuscript, across all sections (Abstract, Introduction, theoretical development, proof sketches, applications, conclusion, and appendix) and across all stages from initial drafting to final revision. The tools were used for boilerplate generation, prose and notation refinement, \LaTeX{}/structure cleanup, translation of informal proof ideas into candidate formal artifacts (Lean/\LaTeX{}), and repeated adversarial reviewer-style critique passes to identify blind spots and clarity gaps.

The author retained full intellectual and editorial control, including problem selection, theorem statements, assumptions, novelty framing, acceptance criteria, and final inclusion/exclusion decisions. No technical claim was accepted solely from AI output. Formal claims reported as machine-verified were admitted only after Lean verification (no \texttt{sorry} in cited modules) and direct author review; Lean was used as an integrity gate for responsible AI-assisted research. The author is solely responsible for all statements, citations, and conclusions.

\subsection*{Data and Code Availability}

The Lean~4 formalization, supplementary proofs, manuscript PDF, and source bundle supporting the cited theorem handles are archived at Zenodo~\cite{paper2_partial_views_artifact}. No separate experimental dataset is associated with this theoretical study.
